%File: cac_aaai_v3.tex
\documentclass[letterpaper]{article} % DO NOT CHANGE THIS
\usepackage[submission]{aaai2026}     % SWAP to aaai2027 for AAAI-27; keep [submission] for review
\usepackage{times}  % DO NOT CHANGE THIS
\usepackage{helvet} % DO NOT CHANGE THIS
\usepackage{courier} % DO NOT CHANGE THIS
\usepackage[hyphens]{url} % DO NOT CHANGE THIS
\usepackage{graphicx} % DO NOT CHANGE THIS
\urlstyle{rm} % DO NOT CHANGE THIS
\def\UrlFont{\rm} % DO NOT CHANGE THIS
\usepackage{natbib} % DO NOT CHANGE THIS AND DO NOT ADD ANY OPTIONS TO IT
\usepackage{caption} % DO NOT CHANGE THIS AND DO NOT ADD ANY OPTIONS TO IT
\frenchspacing % DO NOT CHANGE THIS
\setlength{\pdfpagewidth}{8.5in} % DO NOT CHANGE THIS
\setlength{\pdfpageheight}{11in} % DO NOT CHANGE THIS

\usepackage{amsmath}
\usepackage{amssymb}
\usepackage{booktabs}
\usepackage{multirow}
\usepackage{algorithm}
\usepackage{algorithmic}

\setcounter{secnumdepth}{0}

% /Title required by AAAI 2026 guidelines. (Paper title is not identifying.)
\pdfinfo{
/Title (Before You Trust Disagreement-Based Routing: A Router Acceptance Protocol and Its Validation Across Scale, Modality, and Domain)
/TemplateVersion (2026.1)
}

\title{Before You Trust Disagreement-Based Routing:\\
A Router Acceptance Protocol and Its Validation Across Scale, Modality, and Domain}

% ANONYMIZED FOR DOUBLE-BLIND REVIEW.
% Do NOT restore name/affiliation/email until the camera-ready (non-submission) build.
\author{
    Anonymous Submission
}
\affiliations{
    Paper submitted for double-blind review
}

\begin{document}

\maketitle

\begin{abstract}
A popular heuristic holds that when a small, cheap ensemble of models disagrees
on an input, that input is genuinely hard, so disagreement can route annotation,
escalating contested items to an expensive frontier model or human. We argue this
heuristic is trusted too readily, and turn the implicit cautionary advice of prior
work into a concrete artifact: a \textbf{router acceptance protocol} that
practitioners run \emph{before} deploying disagreement-based routing. The protocol
prescribes four pre-deployment checks --- (i) selective-risk utility (AURC), not
just correlation; (ii) a within-family capacity sweep; (iii) an
error-\emph{independence} test across ensemble members; and (iv) a \emph{measured}
frontier-accuracy tier rather than an assumed one --- each with a pass/fail rule.
We validate the protocol on human label distributions and human grading across
four datasets, two modalities, two tasks, and a capacity sweep from $0.8$B to a
frontier model. Three findings justify the checks. First, routing utility can
\emph{decouple} from correlation: a calibrated cascade beats its single-model
baseline at $4.4\times$ lower cost even when the disagreement--difficulty
correlation is modest. Second, correlation rises with capacity but
non-monotonically, so two-point comparisons are unsafe; AURC is the more
interpretable metric. Third, small dense models can exhibit strongly correlated errors
($\phi{=}0.78$--$0.83$), weakening the independence assumption behind
disagreement-based routing; genuine independence appears only at a
mixture-of-experts model in the language setting. A frontier
control shows its uncertainty tracks human disagreement ($r{=}0.515$) where
small-model disagreement does not, confirmed on two human-graded datasets. We
release the protocol as a reusable acceptance gate.
\end{abstract}

\begin{links}
\link{Code}{\url{https://anonymous.4open.science/r/router-acceptance-protocol-71E4/}}
\end{links}

\section{Introduction}
Supervised systems are bottlenecked by labeled data, and the labels themselves
are imperfect: even large benchmarks built with intensive crowd effort
\citep{deng2009imagenet} contain errors, and careful re-annotation reveals
pervasive mistakes and genuine ambiguity in standard test sets
\citep{northcutt2021pervasive, beyer2020imagenet}.
A now-common practice is to pre-label new data with a capable model and reserve
scarce human (or expensive frontier-model) effort for the items that need it.
This reframes annotation as a \emph{routing} problem: under a fixed budget, decide
which items a cheap labeler can be trusted to finalize and which must escalate.

The hard part is the decision rule. A single model's own confidence is a weak
proxy \citep{hendrycks2017baseline, guo2017calibration}, and it cannot separate
``easy and certain'' from ``hard but overconfident.'' A widely held intuition
offers an alternative: \emph{ensemble disagreement}. When several diverse models
disagree, the input is plausibly ambiguous
\citep{lakshminarayanan2017simple, gal2016dropout}, so disagreement might rank
items by difficulty and trigger escalation. The intuition is attractive precisely
because it is cheap: it asks small models a question they can answer in one
forward pass, and it needs no labels to deploy.

That cheapness is also its danger. Disagreement-based routing is easy to stand up,
easy to demo on a small sample, and easy to \emph{trust} on the strength of a
single encouraging correlation --- and, as we show, that trust is frequently
misplaced. The contribution of this paper is to make the trust decision explicit.
We promote the implicit ``check before you rely on it'' advice of prior work to a
\textbf{first-class artifact}: a \emph{router acceptance protocol} a practitioner
runs before deploying disagreement-based routing, with concrete checks and
pass/fail rules. We then validate the protocol by running it ourselves across
model scale, two modalities, and two task types, and showing that each check
catches a real failure that a naive correlation test would have missed.

The intuition is testable because the target is observable. CIFAR-10H
\citep{peterson2019human} and ChaosNLI \citep{nie2020chaosnli} each provide
\emph{soft human labels} --- full label distributions from $\sim\!50$--$100$
annotators per item --- whose entropy quantifies per-item human disagreement; two
human-graded short-answer and code-grading datasets provide per-item human score
distributions for an open-ended second domain. We can therefore measure whether
model disagreement tracks human difficulty, and separately whether the resulting
routing is useful, across modality, scale, and task.

\paragraph{The router acceptance protocol (lead contribution).}
Before trusting disagreement-based routing, run four checks:
\begin{enumerate}\itemsep2pt
\item[\textbf{A1}] \textbf{Utility, not correlation.} Report selective risk
(AURC) and selective accuracy under a budget, not only the
disagreement--difficulty correlation. \emph{Pass} if the calibrated cascade
matches or beats its single-model baseline at the target budget; a good
correlation with no downstream win is a fail, and (as we show) a modest
correlation with a real win is a pass.
\item[\textbf{A2}] \textbf{Capacity sweep, not two points.} Trace the metric
across a within-family capacity range. \emph{Fail} the ``it's a property of
small models'' claim if it rests on a two-point contrast; the relationship is
non-monotonic and a single pair can mislead in either direction.
\item[\textbf{A3}] \textbf{Error independence, not just disagreement.} Measure
the pairwise correlation of ensemble members' \emph{errors}. \emph{Fail} the
ensemble premise if member errors are correlated --- disagreement then carries no
independent signal and ensembling cannot help.
\item[\textbf{A4}] \textbf{Measured frontier, not assumed.} Measure the escalation
tier's accuracy on the actual hard items rather than assuming a ceiling (e.g.\
$0.95$). \emph{Fail} any cost/accuracy projection built on an unmeasured frontier
number; measured frontier accuracy is domain-specific and often far lower.
\end{enumerate}

\paragraph{Validation findings (supporting contributions).}
\begin{enumerate}\itemsep2pt
\item \textbf{Utility decouples from correlation (validates A1).} On
$10{,}000$-image CIFAR-10H a calibrated cascade exceeds its best single model
($0.965$ vs.\ $0.940$) at $4.4\times$ lower frontier cost despite a modest
correlation ($r{=}0.32$): a router need not rank all items, only isolate the ones
it would mislabel. A disagreement score can support effective routing even when
its global correlation with human difficulty is modest.
\item \textbf{Capacity is a curve, and AURC is the better axis (validates A2).}
A within-family sweep ($0.8$B$\to$frontier) shows disagreement-tracking rises with
scale but non-monotonically (a $4$B checkpoint is a marked outlier); selective
risk (AURC) decreases with scale more cleanly than correlation and is the more
interpretable acceptance metric.
\item \textbf{Small dense models err \emph{together} (validates A3; reverses a
prior claim).} The belief that small models are confidently wrong in
\emph{uncorrelated} ways is false for small dense models: pairwise error
correlation is $\phi{=}0.78$--$0.83$. Error independence appears only at a
mixture-of-experts model. We therefore retract the uncorrelated-error claim and
report the correlated-error finding with its evidence.
\item \textbf{A measured frontier and a second domain (validates A4).} A frontier
control on ChaosNLI tracks human disagreement at $r{=}0.515$ where small-model
disagreement gives $r{\approx}0.13$; its measured accuracy on hard items
($0.673$) is far below an illustrative $0.95$ ceiling that cost projections
might otherwise assume. Two
human-graded datasets (short-answer and code grading) reproduce the
disagreement-does-not-track-difficulty null in an open-ended domain.
\end{enumerate}

\section{Related Work}

\paragraph{Human label uncertainty.}
\citet{peterson2019human} collected CIFAR-10H and showed soft human labels improve
robustness; \citet{nie2020chaosnli} collected ChaosNLI (${\sim}100$ annotations per
NLI example) and showed that strong models only partially recover human label
\emph{distributions}, with accuracy collapsing on high-disagreement items.
Re-annotation studies find ``ground-truth'' labels are often wrong or ambiguous
\citep{northcutt2021pervasive, beyer2020imagenet}. We use the label-distribution
entropy of both as our measure of per-item human disagreement, and extend to human
\emph{score} distributions for open-ended grading.

\paragraph{Ensembles, calibration, uncertainty.}
Deep ensembles \citep{lakshminarayanan2017simple}, MC-dropout
\citep{gal2016dropout}, and member disagreement \citep{kendall2017uncertainties}
estimate epistemic uncertainty; calibration is handled by Platt/temperature
scaling \citep{guo2017calibration, platt1999probabilistic}, isotonic regression
\citep{zadrozny2002transforming}, and ECE \citep{naeini2015obtaining}. A standing
assumption behind ensembling is that members make \emph{independent} errors
\citep{dietterich2000ensemble}; our
A3 check tests that assumption directly and finds it violated for small dense
models. Throughout this paper, \emph{disagreement} (unless qualified) means
\emph{distributional} disagreement among member posteriors --- the mean pairwise
Jensen--Shannon divergence defined in Table~\ref{tab:notation} --- which we keep
distinct from label-vote disagreement and from error disagreement (the $\phi$
correlation of error indicators in A3). Our objective is not posterior accuracy
but the \emph{ranking} of items by human-perceived difficulty, and crucially
\emph{whether} this distributional disagreement supports it.

\paragraph{Selective prediction and cost-aware cascades.}
Selective classifiers \citep{geifman2017selective, hendrycks2017baseline} and
conformal methods \citep{angelopoulos2021gentle} abstain under uncertainty;
cascades \citep{viola2001rapid} and FrugalGPT \citep{chen2023frugalgpt} invoke
expensive models only when cheap ones seem insufficient; LLM-as-judge
\citep{zheng2023judging} gates computation on model outputs. We adopt selective
risk (AURC) as our headline acceptance metric and use a disagreement-triggered
cascade as the worked system the protocol validates.

\paragraph{Weak supervision and active learning.}
Programmatic weak supervision \citep{ratner2017snorkel} and active learning
\citep{settles2009active} prioritize annotation effort; our protocol is
complementary, auditing when a disagreement-based difficulty estimator may be
trusted for triage.

\paragraph{Models.}
For the within-family capacity sweep we use the Qwen3.5 family (natively
multimodal; $0.8$B--$122$B) so that scale varies within a single backbone, removing
the architecture confound that a cross-family ensemble would introduce. The
cascade system additionally uses open VLMs --- Qwen2-VL \citep{wang2024qwen2vl},
InternVL \citep{chen2024internvl}, Phi-3.5-Vision \citep{abdin2024phi3} --- and
open LLMs (Qwen2.5, Gemma-2). Our contribution is not a new model but a protocol
for deciding when small ensembles of such models can be trusted as uncertainty
sensors.

\paragraph{Recent routing, VLM confidence, and judge work (2024--2026).}
LLM routing has matured into learned cost-quality trade-offs: RouteLLM
\citep{ong2024routellm} trains a router from preference data to send easy queries
to a weak model and hard ones to a strong model. Our setting differs in the
trigger --- cross-model \emph{disagreement} against a human-difficulty target, not a
learned strong-vs-weak win predictor --- and in the target, human label entropy
rather than response preference. Recent work documents that VLMs are
systematically miscalibrated and that verbalized confidence often misaligns with
accuracy \citep{groot2024verbalized}, and that sampling-based uncertainty only
partially recovers reliability \citep{zhang2024vluncertainty}. These findings are
consistent with our small-model results and motivate the measured-frontier check
(A4).

\section{Problem Setup}
For each item $x$ over $K$ classes (or a bounded score range), the dataset provides
a human label/score distribution $p_h(\cdot\mid x)$ (empirically observed from
annotators) with entropy
\begin{equation}
H_{\text{human}}(x) = -\sum_{k=1}^{K} p_h(k\mid x)\log p_h(k\mid x),
\end{equation}
which we take as our measure of per-item \emph{human disagreement}. A router
assigns each item an action
$a\in\{\textsf{skip},\textsf{consensus},\textsf{frontier}\}$ with costs
$c_{\textsf{skip}}\le c_{\textsf{consensus}}\ll c_{\textsf{frontier}}$. With
$\rho_a$ the fraction taking action $a$,
\begin{equation}
C = N\big(c_{\text{ens}} + \rho_{\textsf{frontier}}\,c_{\textsf{frontier}}\big),
\label{eq:cost}
\end{equation}
so the saving over a frontier-only policy is $C_F/C\approx
1/\rho_{\textsf{frontier}}$. The goal is to escalate exactly the items the
ensemble would get wrong. Whether disagreement approximates that target ---
and, crucially, \emph{how to decide} whether it does before deploying --- is the
question this paper makes explicit. Table~\ref{tab:notation} fixes the notation and
the three distinct senses of ``disagreement'' used below.

\begin{table}[t]
\centering
\small
\begin{tabular}{lp{0.62\columnwidth}}
\toprule
\textbf{Symbol / term} & \textbf{Definition} \\
\midrule
$p_h(\cdot\mid x)$ & Observed human label/score distribution for item $x$ over $K$ classes (or a bounded score range). \\
$H_{\text{human}}(x)$ & Human-label entropy, $-\sum_k p_h(k\mid x)\log p_h(k\mid x)$; our measure of per-item \emph{human disagreement} (observed, not assumed difficulty). \\
$p_m(\cdot\mid x)$ & Predictive distribution of ensemble member $m$ on item $x$. \\
\textbf{Disagreement} & Mean pairwise Jensen--Shannon divergence among the $M$ member distributions, $\mathrm{JSD}(x)=\binom{M}{2}^{-1}\sum_{i<j}\mathrm{JS}(p_i\Vert p_j)$ (natural log). Unless stated otherwise, ``disagreement'' refers to this. \\
$\mathrm{JS}(p\Vert q)$ & $\tfrac12 D_{\mathrm{KL}}(p\Vert m)+\tfrac12 D_{\mathrm{KL}}(q\Vert m)$, with $m=\tfrac12(p+q)$. \\
\textbf{Correlation} & Pearson $r$ (and Spearman $\rho$) between per-item disagreement and $H_{\text{human}}$ over the $N$ items; the relationship the protocol declines to use as the \emph{sole} acceptance test. \\
\textbf{Error indep.} & Pairwise $\phi$ coefficient between members' binary per-item \emph{error} indicators on shared items; low $\phi$ $=$ idiosyncratic errors. \\
$s(x)$ & Escalation signal $s(x)=\alpha\,\widetilde{\mathrm{JSD}}(x)+\beta\,(1{-}\mathrm{MTA}(x))$; $\widetilde{\cdot}$ is min--max normalized, $\alpha,\beta$ set by mutual information with $H_{\text{human}}$ ($\beta{=}0$ recovers JSD-only). \\
$\mathrm{MTA}$ & Mean textual-rationale agreement among members (sentence-embedding cosine similarity); $1{-}\mathrm{MTA}$ is rationale dissimilarity. \\
$\mathrm{ECE}$ & Expected calibration error of the calibrated $P(\text{hard})$, $10$ equal-width bins: $\sum_b \tfrac{|B_b|}{N}\,|\mathrm{acc}(B_b)-\mathrm{conf}(B_b)|$. \\
\textbf{False-easy rate} & Fraction of truly-hard items the pre-router admits as easy: $\#\{\text{skipped}\wedge\text{hard}\}/\#\{\text{hard}\}$. \\
\textbf{Sel.\ acc.\ @$c$} & Accuracy on the fraction $c$ of items with lowest escalation score (the remainder escalate). \\
$\mathrm{AURC}$ & Area under the risk--coverage curve (lower better); risk $=$ error rate at each coverage of the escalation score. \\
$\tau,\ \rho_{\textsf{frontier}}$ & Escalation threshold; fraction escalated to the frontier tier under budget. \\
\bottomrule
\end{tabular}
\caption{Notation. ``Disagreement'' denotes mean pairwise JSD unless otherwise
specified; ``correlation'' denotes Pearson/Spearman of disagreement against
$H_{\text{human}}$.}
\label{tab:notation}
\end{table}

\section{The Router Acceptance Protocol}
The protocol is a pre-deployment gate. Each check produces a pass/fail and a
reusable artifact (a number with a confidence interval, a curve, or a matrix).
Table~\ref{tab:protocol} summarizes the four checks as a failure-mode $\to$
diagnostic $\to$ pass/fail rule; the criteria are specified independently of our
results, and the parenthetical numbers are what we observed when we ran the
protocol ourselves.

\begin{table*}[t]
\centering
\small
\begin{tabular}{p{0.03\textwidth}p{0.22\textwidth}p{0.27\textwidth}p{0.36\textwidth}}
\toprule
& \textbf{Failure mode it catches} & \textbf{Diagnostic check (artifact)} & \textbf{Pass / fail criterion} \\
\midrule
\textbf{A1} & Trusting a correlation that yields no routing benefit (or rejecting a useful router for low $r$). & Calibrated cascade vs.\ single-model baseline; AURC and selective accuracy at the budget (risk--coverage curve). & \emph{Pass} if the calibrated cascade matches or beats the best single model at target coverage (lower AURC, $\ge$ selective accuracy). \emph{Fail} if a good $r$ produces no downstream win. \\
\addlinespace
\textbf{A2} & Attributing behavior to ``small models'' from a two-point contrast a non-monotone curve would contradict. & AURC and $r$ traced across a within-family capacity range (a curve, not a pair); outliers flagged. & \emph{Pass} only with a characterized trend across $\ge\!4$ capacity points. \emph{Fail} any claim resting on two points; flag any point whose AURC residual exceeds ${\sim}2\times$ the next largest. \\
\addlinespace
\textbf{A3} & Assuming ensemble members err independently when they fail on the same items. & Pairwise $\phi$ of member error indicators on shared items (error-correlation matrix). & \emph{Pass} if mean pairwise $\phi$ is low (idiosyncratic errors). \emph{Fail} the ensemble premise if errors are strongly correlated (here $\phi{=}0.78$--$0.83$ for small dense members). \\
\addlinespace
\textbf{A4} & Cost/accuracy projections built on an assumed frontier ceiling rather than a measured one. & Frontier tier's measured accuracy on the actual hard subset, with CI. & \emph{Pass} only if projections use the \emph{measured} frontier accuracy for the domain. \emph{Fail} any projection substituting an assumed ceiling (here $0.95$ assumed vs.\ $0.673$ measured). \\
\bottomrule
\end{tabular}
\caption{The router acceptance protocol as a reusable gate. Each check targets a
specific failure mode, produces a reusable artifact, and has an operational
pass/fail rule. Criteria are specified independently of our experiments;
parenthetical values are what we observed when running the protocol.}
\label{tab:protocol}
\end{table*}

\paragraph{A1 --- Utility over correlation.}
Compute the area under the risk--coverage curve (AURC) and selective accuracy at
the target coverage, using a calibrated escalation score. \emph{Decision:} accept
only if the calibrated cascade matches or beats its single-model baseline at the
budget. Correlation alone is neither necessary nor sufficient (Sec.~Results, A1).

\paragraph{A2 --- Capacity sweep.}
Trace AURC and correlation across a within-family capacity range rather than a
single small-vs-large pair. \emph{Decision:} do not attribute behavior to ``small
models'' from two points; require a monotone or clearly-characterized trend, and
flag outliers explicitly.

\paragraph{A3 --- Error independence.}
For the proposed ensemble members, compute the pairwise correlation
($\phi$ coefficient) of their per-item \emph{error indicators} on shared items.
\emph{Decision:} reject the ensemble premise if member errors are strongly
correlated; disagreement then adds no independent information over a single member.

\paragraph{A4 --- Measured frontier.}
Measure the escalation tier's accuracy on the actual hard subset. \emph{Decision:}
reject any cost/accuracy projection that substitutes an assumed ceiling for a
measured one; report the measured number and its domain.

\paragraph{A concrete failure the protocol catches.}
Consider a practitioner evaluating a disagreement-based router on ChaosNLI.
They observe a modest positive disagreement--difficulty correlation and conclude
that disagreement is a useful proxy for difficulty; a correlation-only analysis
would accept the router. The protocol, however, exposes two hidden failures:
(i)~the router delivers little selective-risk benefit~(A1), and (ii)~the ensemble
members make highly correlated errors ($\phi \approx 0.80$), violating the
independence assumption that motivates disagreement-based routing~(A3). The
protocol therefore \emph{rejects} a router that a correlation-only evaluation
would have accepted.

\paragraph{The disagreement-based routing pipeline under test.}
We instantiate routing with a disagreement-based cascade: each item is run through
$M$ small models; a distributional disagreement score (Jensen--Shannon
divergence \citep{lin1991divergence} among per-model posteriors) is optionally
fused with a textual rationale-dissimilarity signal
\citep{reimers2019sentencebert}, calibrated to an escalation probability
\citep{platt1999probabilistic, zadrozny2002transforming, naeini2015obtaining}, and
thresholded under a frontier budget to drive a skip/consensus/frontier cascade
(Algorithm~\ref{alg:cac}). This pipeline is the worked example the protocol audits,
not the paper's contribution.

\paragraph{Rationale agreement (MTA).}
Beyond distributional disagreement, each member emits a short textual rationale for
its answer. We encode rationales with a sentence-embedding model
\citep{reimers2019sentencebert} and define \emph{mean textual agreement} (MTA) as
the average pairwise cosine similarity among member rationales on an item; high MTA
means members justify their answers similarly. The complement $1{-}\mathrm{MTA}$ is
a second, text-level disagreement channel that is distinct from the
distribution-level JSD. The escalation signal fuses the two,
$s(x)=\alpha\,\widetilde{\mathrm{JSD}}(x)+\beta\,(1{-}\mathrm{MTA}(x))$, where the
weights $\alpha,\beta$ are set by each channel's mutual information with
$H_{\text{human}}$; when rationales add no information beyond the posteriors, the
fit drives $\beta\!\to\!0$ and the signal reduces to JSD alone.

\begin{algorithm}[t]
\caption{Disagreement-based routing for an item $x$}
\label{alg:cac}
\begin{algorithmic}[1]
\STATE Run ensemble $\{p_m, r_m\}_{m=1}^{M}$ on $x$
\STATE Compute disagreement score $s(x)$ (JSD, optionally fused with rationale dissimilarity)
\STATE $e \gets g(s(x))$ \COMMENT{calibrated escalation probability}
\IF{pre-router marks $x$ high-confidence consensus}
   \STATE \textbf{skip}: accept ensemble-majority label
\ELSIF{$e \le \tau$}
   \STATE \textbf{consensus}: accept confidence-weighted label
\ELSE
   \STATE \textbf{frontier}: escalate $x$
\ENDIF
\end{algorithmic}
\end{algorithm}

\section{Experimental Setup}

\paragraph{Datasets and target.}
\emph{Vision:} CIFAR-10H \citep{peterson2019human}, the $10{,}000$ CIFAR-10
\citep{krizhevsky2009learning} test images with soft human labels; human-argmax
agrees with the original labels at $0.992$, so hard items are a small tail. For the
capacity sweep we use a frozen $450$-item subset stratified by human entropy that
oversamples the ambiguous tail (so the disagreement question is answerable on a
mostly-zero-entropy benchmark).
\emph{Language:} ChaosNLI \citep{nie2020chaosnli}, the SNLI$+$MNLI subset
($N{=}3{,}113$, three classes, ${\sim}100$ annotators each; mean human entropy
$0.94$ bits); the capacity sweep uses a matched $450$-item stratified subset so
every model and the frontier control are scored on identical items.
\emph{Open-ended grading (second domain):} two human-graded datasets --- a
short-answer grading set ($N{=}2{,}254$) and a programming-assignment grading set
($N{=}998$ after filtering ungradeable submissions), each with multiple human
graders per item --- provide per-item human score distributions.

\paragraph{Capacity sweep.}
Qwen3.5 at $0.8$B, $2$B, $4$B, $9$B, $27$B (dense) and $122$B
(mixture-of-experts; $\sim\!10$B active). Each model is asked for a full
class-probability distribution per item in one structured call, and we take its
entropy as the model's uncertainty. Hosted endpoints serve
$0.8$B/$27$B/$122$B; $2$B/$4$B/$9$B run locally via vLLM. The frontier control is
Claude Sonnet~4.6.

\paragraph{Cascade ensembles.}
\emph{Vision 2B:} Qwen2-VL-2B + InternVL2-2B (the system that demonstrates A1).
\emph{Language:} Qwen2.5-3B-Instruct + Gemma-2-2B-it for the prior-generation
ensemble comparison.

\paragraph{Metrics and decisions.}
Pearson/Spearman of model uncertainty vs.\ $H_{\text{human}}$; \textbf{AURC}
(area under the risk--coverage curve, lower is better) and selective accuracy at
$50/70/90\%$ coverage as the headline utility metrics; AUROC for ranking hard
items; pairwise error-correlation ($\phi$) for A3; single-model vs.\ cascade
accuracy, escalated fraction, and cost reduction (Eq.~\ref{eq:cost}) for the
system. We report bootstrap $95\%$ CIs.

\paragraph{Implementation and compute.}
vLLM constrained decoding, fixed seed, offline weights; cross-encoder stack pinned.
The capacity-sweep models run as resumable, content-cached jobs. Experiments were
conducted on one compute node with two NVIDIA A100 GPUs (driver version
595.71.05, CUDA 13.2), 20 Intel Xeon Platinum 8468 CPUs, and 234\,GB RAM. Code and
configs to be released via the anonymized repository linked above.

\section{Results}

\subsection{A1 --- Routing utility can decouple from correlation (CIFAR-10H)}
At $N{=}10{,}000$ the 2B cascade reaches $0.965$, above its best single model
($0.940$), escalating $22.8\%$ for a $4.4\times$ cost reduction at
$\mathrm{ECE}{=}0.015$ --- despite a modest disagreement--entropy correlation
($r{=}0.32$). A calibrated threshold fit on enough data isolates the
sub-population the ensemble would mislabel even from a weakly-correlated score; the
pre-router's false-easy rate falls from ${\approx}37\%$ at $N{=}500$ to $1.2\%$ at
$N{=}10{,}000$, the mechanism behind the gain. Utility and global correlation are
empirically distinct --- which is exactly why A1 requires the utility measurement,
not the correlation, as the acceptance test.

\subsection{A2 --- Capacity is a curve, and AURC is the cleaner axis (ChaosNLI)}
Table~\ref{tab:sweep} reports the within-family sweep on the matched $450$-item
ChaosNLI subset. Disagreement-tracking rises with capacity but \emph{not}
monotonically: the smallest models barely track human disagreement
($r{=}0.138$ at $0.8$B, $0.132$ at $2$B), the signal strengthens through the
mid-range, and the frontier model is clearly best ($r{=}0.471$). Selective risk
tells the cleaner story: AURC falls from $0.465$ ($0.8$B) toward $0.178$ (frontier)
and is easier to interpret than $r$ against the zero-inflated entropy target,
which is why the protocol makes AURC the headline (A1/A2). A $4$B checkpoint is a
marked positive outlier (AURC $0.230$, well below the dense trend; see
Figure~\ref{fig:curve}); we treat it as a checkpoint-specific anomaly and do not
build claims on it, but it is precisely the kind of non-monotonicity that a
two-point ``small vs.\ large'' comparison would misread.\footnote{The $4$B point's
AURC residual from a log-linear fit on the other dense models is $-0.18$, roughly
twice the next-largest residual; the upward trend with scale is robust to
excluding it, but the magnitude of the deviation is why we flag rather than smooth
it.} Per-model selective accuracy at $50/70/90\%$ coverage is reported in
Table~\ref{tab:sweep}.

\begin{table}[t]
\centering
\small
\resizebox{\columnwidth}{!}{%
\begin{tabular}{lcccc}
\toprule
\textbf{Model} & $r$ & \textbf{AURC}$\downarrow$ & \textbf{sel@50} & \textbf{sel@70} \\
\midrule
Qwen3.5-0.8B            & $0.138$ & $0.465$ & $0.529$ & $0.486$ \\
Qwen3.5-2B             & $0.132$ & $0.424$ & $0.578$ & $0.606$ \\
Qwen3.5-4B$^{\ast}$    & $0.377$ & $\mathbf{0.230}$ & $0.787$ & $0.686$ \\
Qwen3.5-9B             & $0.304$ & $0.306$ & $0.618$ & $0.584$ \\
Qwen3.5-27B            & $0.271$ & $0.438$ & $0.520$ & $0.476$ \\
Qwen3.5-122B (MoE)$^{\dagger}$ & $0.268$ & $0.387$ & $0.596$ & $0.495$ \\
\midrule
Frontier (Sonnet 4.6)  & $\mathbf{0.471}$ & $\mathbf{0.178}$ & --- & --- \\
\bottomrule
\end{tabular}%
}
\caption{Within-family capacity sweep (ChaosNLI, matched $450$-item subset).
Disagreement-tracking ($r$) and selective risk (AURC, lower better) improve with
scale but non-monotonically; AURC is the cleaner trend and the protocol's headline
metric. $^{\ast}$$4$B is a positive outlier (see text and Fig.~\ref{fig:curve}).
$^{\dagger}$The $122$B is a mixture-of-experts model ($\sim\!10$B active of $122$B
total); it is reported separately and not placed on the dense-capacity trend.
Selective-accuracy entries for the frontier are not applicable: the frontier is the
escalation \emph{target} and answers all routed items, so it makes no selective
decision of its own.}
\label{tab:sweep}
\end{table}

\begin{figure}[t]
\centering
\includegraphics[width=\columnwidth]{fig_4b_outlier.png}
\caption{Selective risk (AURC, lower is better) vs.\ dense model capacity on
ChaosNLI. AURC improves with scale toward the frontier; the $4$B checkpoint sits
well below the dense trend (right panel: residual roughly $2\times$ the next
largest), and the $122$B mixture-of-experts model is shown as a separate marker
rather than on the dense line. This is the non-monotonicity that motivates the
capacity-sweep check (A2): a two-point comparison would misread it.}
\label{fig:curve}
\end{figure}

\subsection{A3 --- Small dense models err \emph{together}, not independently
(ChaosNLI)}
The ensembling premise is that members make independent errors, so disagreement
flags genuine difficulty. We test this directly by correlating per-item error
indicators across the swept models on the shared $450$ items
(Figure~\ref{fig:errcorr}). The result \textbf{reverses} the common
``confidently wrong in uncorrelated ways'' intuition for small dense models: their
errors are strongly correlated --- $\phi{=}0.815$ ($0.8$B vs.\ $2$B), $0.779$
($0.8$B vs.\ $27$B), $0.828$ ($2$B vs.\ $27$B) --- i.e.\ they fail on the
\emph{same} items, so an ensemble of them carries little independent signal and
cannot separate hard from easy by disagreement. Error independence appears only at
the $122$B mixture-of-experts model, whose error vector is near-uncorrelated with
every dense model ($\phi{\in}[-0.05, 0.14]$). We therefore retract the
uncorrelated-error claim of the prior version of this work and report the
correlated-error finding as evidence for the A3 check: measuring error
independence, not merely observing disagreement, is what reveals that small dense
ensembles cannot deliver the premise.

\begin{figure}[t]
\centering
\includegraphics[width=\columnwidth]{fig_error_correlation.png}
\caption{Pairwise error correlation ($\phi$) among swept models on shared ChaosNLI
items (low $=$ idiosyncratic errors). Small dense models form a strongly-correlated
block ($0.78$--$0.83$): they err on the same items, so disagreement among them is
not independent evidence of difficulty. Only the $122$B mixture-of-experts model
shows near-independent errors. This violates the ensembling premise and is the
basis for acceptance check A3.}
\label{fig:errcorr}
\end{figure}

\subsection{A4 --- A measured frontier, not an assumed one (ChaosNLI)}
Cost/accuracy projections for cascades commonly assume a frontier ceiling near
$0.95$ on escalated items. We measure it. On the stratified ChaosNLI subset the
frontier model (Claude Sonnet~4.6), queried for a full class-probability
distribution per item, reaches accuracy $0.673$ against the human majority --- far
below the assumed $0.95$, and only modestly above chance on the hardest tercile.
Critically, its \emph{uncertainty} tracks human disagreement at $r{=}0.515$
($95\%$ CI $[0.446,0.578]$), an order of magnitude above the small models'
$r{\approx}0.13$; across entropy terciles its mean entropy rises monotonically as
accuracy falls, so it grows appropriately less certain on exactly the items humans
find ambiguous (Table~\ref{tab:frontier}). Two consequences follow. First, the
small-model failure is a property of capacity, not of the task --- a capable model
recovers the disagreement--difficulty relationship on the same items. Second, any
cost projection must use the \emph{measured} frontier accuracy for the target
domain ($0.673$ here), not a generic ceiling; substituting $0.95$ would overstate
the cascade's achievable accuracy.

\begin{table}[t]
\centering
\small
\resizebox{\columnwidth}{!}{%
\begin{tabular}{lccc}
\toprule
\textbf{System (ChaosNLI 450-subset)} & \textbf{Acc} & \textbf{corr(unc.,$H$)} & \textbf{95\% CI} \\
\midrule
Small models (Qwen3.5 $\le 2$B) & ${\sim}0.49$ & ${\sim}0.13$ & --- \\
\textbf{Frontier (Sonnet 4.6)} & $\mathbf{0.673}$ & $\mathbf{0.515}$ & $\mathbf{[0.446,0.578]}$ \\
\bottomrule
\end{tabular}%
}
\caption{Measured frontier tier (A4). A frontier model's stated uncertainty tracks
human label entropy where small-model uncertainty does not; its \emph{measured}
accuracy on hard items ($0.673$) is far below an illustrative $0.95$ ceiling
that cost projections might otherwise assume. This attributes the small-model failure to capacity, not the
task, and supplies the measured number A4 requires.}
\label{tab:frontier}
\end{table}

\subsection{Second domain: human-graded short-answer and code grading}
To test whether the disagreement-does-not-track-difficulty result generalizes
beyond classification, we evaluate two open-ended grading datasets with multiple
human graders per item, using inter-model score disagreement as the routing signal
and human grader spread as difficulty (Table~\ref{tab:grading}). In both, model
disagreement barely tracks human disagreement: short-answer grading gives
$r{=}0.090$ ($\rho{=}0.182$, $N{=}2{,}254$) and code grading gives $r{\approx}0$
($\rho{=}0.034$, $N{=}998$). Selecting items by disagreement yields only a small
lift over random selection ($+1$--$3\%$ on average). One honest caveat governs the
strength of this conclusion: the grading models' own accuracy within tolerance of
the adjudicated score is low (short-answer $0.376$; code $0.215$), so part of the
null could reflect a weak grader rather than a property of disagreement; we report
both the null and this limitation. The two datasets nonetheless reproduce the same
qualitative result in a second domain, which is what the protocol's cross-domain
re-validation prescribes.

\begin{table}[t]
\centering
\small
\resizebox{\columnwidth}{!}{%
\begin{tabular}{lcccc}
\toprule
\textbf{Grading dataset ($N$)} & $r$ & $\rho$ & \textbf{lift} & \textbf{grader acc} \\
\midrule
Short-answer ($2{,}254$) & $0.090$ & $0.182$ & ${+}1.4\%$ & $0.376$ \\
Code grading ($998$)     & $-0.005$ & $0.034$ & ${+}1.2\%$ & $0.215$ \\
\bottomrule
\end{tabular}%
}
\caption{Open-ended grading (second domain). Model disagreement does not track
human grading disagreement; selecting by disagreement gives only a small mean lift
over random. ``lift'' is the mean accuracy gain over random selection across
coverage; ``grader acc'' is model accuracy within tolerance of the adjudicated
score (a limitation, see text).}
\label{tab:grading}
\end{table}

\subsection{Vision capacity sweep generalizes A2, with a modality-specific
A3 result (CIFAR-10H)}
We repeat the within-family capacity sweep on the multimodal CIFAR-10H track
(Qwen3.5 $0.8$B--$27$B; $2$B/$4$B/$9$B served locally via vLLM) on the same
$450$-item stratified subset used elsewhere, to test whether the language-track
findings (A2 capacity curve, A3 error correlation) hold across modality. The
$122$B point is excluded: the hosted endpoint for this model is
text-only and rejects image inputs, so it cannot be placed on the vision curve.
For Qwen3.5-27B, $224/450$ items returned a parseable distribution (the remainder
exhausted the token budget on reasoning before emitting an answer); we report its
metrics over these $224$ items and flag the smaller sample.

\begin{table}[t]
\centering
\small
\begin{tabular}{lcccc}
\toprule
\textbf{Model} & $r$ & \textbf{AURC}$\downarrow$ & \textbf{AUROC} & \textbf{sel@70} \\
\midrule
Qwen3.5-0.8B & $0.137$ & $0.331$ & $0.580$ & $0.692$ \\
Qwen3.5-2B   & $0.376$ & $0.137$ & $0.682$ & $0.876$ \\
Qwen3.5-4B   & $0.447$ & $0.137$ & $0.781$ & $0.829$ \\
Qwen3.5-9B   & $\mathbf{0.529}$ & $\mathbf{0.098}$ & $\mathbf{0.804}$ & $0.857$ \\
Qwen3.5-27B$^{\ddagger}$ & $\mathbf{0.547}$ & $0.206$ & $0.781$ & $0.651$ \\
\bottomrule
\end{tabular}
\caption{Within-family capacity sweep on CIFAR-10H (matched $450$-item subset; the
$122$B endpoint is text-only and excluded). Disagreement-tracking ($r$) rises
\emph{monotonically} with scale, $0.137\to0.547$, and AUROC rises from $0.580$ to
a peak of $0.804$ at $9$B; AURC improves correspondingly from $0.331$ to $0.098$.
$^{\ddagger}$$27$B is computed over $224/450$ parseable items (see text); its
AURC uptick relative to $9$B is consistent with this smaller, less-stable sample
rather than a capability regression.}
\label{tab:vision_sweep}
\end{table}

\begin{figure}[t]
\centering
\includegraphics[width=\columnwidth]{fig_capacity_curve_vision.png}
\caption{CIFAR-10H capacity curve: disagreement--$H_{\text{human}}$ correlation
($r$) and AUROC rise with scale, monotonically for $r$. Unlike the ChaosNLI curve
(Fig.~\ref{fig:curve}), the $4$B checkpoint sits \emph{on} the trend here rather
than as a positive outlier --- the $4$B anomaly observed in A2 is therefore
checkpoint- and modality-specific, not a general property of this model size.}
\label{fig:vision_curve}
\end{figure}

\paragraph{A2 generalizes cleanly to vision.} The capacity-tracking relationship
that was non-monotonic on ChaosNLI (Table~\ref{tab:sweep}) is \emph{monotonic} on
CIFAR-10H: $r$ rises at every step from $0.8$B to $27$B
(Figure~\ref{fig:vision_curve}), and AUROC peaks at $9$B ($0.804$) with the $27$B
dip attributable to its reduced parseable sample. Notably, the $4$B checkpoint that
was a marked positive \emph{outlier} on ChaosNLI (Sec.~A2; Fig.~\ref{fig:curve}) is
squarely on-trend here ($r{=}0.447$, between $2$B's $0.376$ and $9$B's $0.529$).
This is itself evidence for A2's prescription: a capacity-sweep finding from one
modality (the $4$B anomaly) does not transfer to another, so claims built on a
single modality --- like claims built on two capacity points --- require
re-validation before being trusted.

\paragraph{A3 does not generalize: vision errors are markedly less correlated.}
Repeating the A3 diagnostic on the five vision points (Figure~\ref{fig:errcorr_vision})
gives a strikingly different picture than ChaosNLI. Pairwise error correlation
among the small dense vision models ranges $\phi{\in}[0.08, 0.48]$ --- the
$0.8$B--$2$B pair is nearly independent ($\phi{=}0.08$), and even the largest
pairwise value ($4$B--$27$B, $\phi{=}0.48$) is below the \emph{smallest} value in
the ChaosNLI small-dense block ($\phi{=}0.78$). The strongly-correlated
``small dense models err together'' block that drove the A3 finding on ChaosNLI
(Sec.~A3; Fig.~\ref{fig:errcorr}) is therefore \textbf{not} a general property of
small dense models: it is, at least partly, a language-task-specific result. We
report this as a second instance of the same methodological point as the $4$B
case above --- an A3 conclusion drawn from one modality should not be assumed to
hold in another --- and leave a mechanistic explanation (e.g.\ whether CIFAR-10's
coarser $10$-way label space simply admits more ways for small models to differ)
to future work.

\begin{figure}[t]
\centering
\includegraphics[width=\columnwidth]{fig_error_correlation_vision.png}
\caption{Pairwise error correlation ($\phi$) among the five CIFAR-10H capacity
points (the $122$B is excluded, as in Table~\ref{tab:vision_sweep}). Compare to
Figure~\ref{fig:errcorr}: vision errors are markedly less correlated across the
board ($\phi{\le}0.48$), so the ChaosNLI A3 finding (small dense models err
together, $\phi{=}0.78$--$0.83$) does not replicate in this modality.}
\label{fig:errcorr_vision}
\end{figure}

\section{Discussion}

\paragraph{Correlation is the wrong acceptance test for a router.}
A naive deployment decision reads a single correlation and trusts it. Our results
show why that fails in both directions: CIFAR at scale has a modest correlation yet
routes well (A1), while small-model NLI has a near-zero correlation \emph{and}
correlated member errors (A3) so it cannot route at all. The protocol replaces the
single-number test with utility (AURC/selective accuracy), a capacity curve,
error-independence, and a measured frontier.

\paragraph{When does disagreement-as-difficulty hold?}
Our results place it in a narrow regime: it requires members whose \emph{errors}
are actually independent (which small dense models are not, A3), enough capacity to
track human ambiguity at all (a curve, not a point, A2), and a calibrated operating
point where utility materializes (A1). The frontier control sharpens the boundary:
language is not intrinsically hostile to the signal --- a capable model tracks human
disagreement on the same items --- but the cost-saving premise is to route with
\emph{cheap} models, and at that price point the signal and the error-independence
both fail. The regime is jointly defined by capacity, modality, and the
member-independence the ensemble premise assumes.

\section{Broader Impact}
Routing contested items toward human or frontier review rather than silently
auto-labeling them can reduce label-error propagation
\citep{northcutt2021pervasive}. The same machinery could be tuned to minimize
oversight for cost; we frame the frontier budget as a quality lever and the
acceptance protocol as a guard against deploying a difficulty signal that does not
exist. A difficulty signal calibrated on one annotator pool may mis-prioritize on
populations whose ambiguity differs, so the calibrator should be re-fit per domain
and re-validated per modality before being trusted.

\section{Limitations and Future Work}
(i)~The capacity sweep elicits a stated distribution per model rather than a
sampled one; sampling at temperature returned near-identical forced labels and
carried no uncertainty, so the stated distribution is the standard
ChaosNLI-style comparison, but a multi-model frontier panel would tighten the A4
attribution. (ii)~The $4$B outlier is flagged, not explained; a calibration- or
thinking-budget analysis of that checkpoint is future work. (iii)~The grading
nulls are partly limited by low grader accuracy (Table~\ref{tab:grading}); a
stronger grader would test whether the disagreement null persists. (iv)~The CIFAR
and prior-generation cascade results use earlier model families
(Qwen2-VL/2.5-class) than the Qwen3.5 capacity sweep; we present them as
complementary evidence across model generations, and a single-generation
end-to-end replication is the cleanest next step. (v)~The vision capacity sweep is
now complete and generalizes A2 cleanly (a monotone curve, with the $4$B checkpoint
on-trend) but not A3 (the small-dense correlated-error block does not replicate on
CIFAR-10H); we read this as evidence that A3 conclusions are domain-specific until
a mechanistic account is established, rather than evidence against the protocol
itself. (vi)~Cost is measured
in frontier-call units; an end-to-end dollar/wall-clock accounting would
strengthen the cost claims.

\section{Conclusion}
The intuition that small-ensemble disagreement tracks human difficulty is trusted
too readily. We turn the implicit ``check first'' advice into a concrete router
acceptance protocol --- utility over correlation, a capacity curve over two points,
error-independence over mere disagreement, and a measured frontier over an assumed
one --- and validate each check against human label and score distributions across
two modalities and two task types. Routing utility can decouple from correlation;
disagreement-tracking is a non-monotone curve in capacity for which AURC is the
cleaner axis; small dense models tend to err \emph{together} rather than independently
in the language setting, a result that does not fully replicate in vision; and a measured frontier tracks human difficulty where
small models do not, at an accuracy far below the assumed ceiling. We recommend
that disagreement-based routing be accepted only after passing these checks, and we
release them as a reusable gate.

\begin{thebibliography}{}

\bibitem[Abdin et~al.(2024)]{abdin2024phi3}
Abdin, M.; et~al. 2024.
\newblock Phi-3 Technical Report.
\newblock \emph{arXiv:2404.14219}.

\bibitem[Angelopoulos and Bates(2021)]{angelopoulos2021gentle}
Angelopoulos, A.~N.; and Bates, S. 2021.
\newblock A Gentle Introduction to Conformal Prediction.
\newblock \emph{arXiv:2107.07511}.

\bibitem[Beyer et~al.(2020)]{beyer2020imagenet}
Beyer, L.; H\'enaff, O.~J.; Kolesnikov, A.; Zhai, X.; and van~den Oord, A. 2020.
\newblock Are We Done with ImageNet?
\newblock \emph{arXiv:2006.07159}.

\bibitem[Chen et~al.(2024)]{chen2024internvl}
Chen, Z.; et~al. 2024.
\newblock InternVL: Scaling Up Vision Foundation Models and Aligning for
Generic Visual-Linguistic Tasks.
\newblock In \emph{CVPR}.

\bibitem[Chen, Zaharia, and Zou(2023)]{chen2023frugalgpt}
Chen, L.; Zaharia, M.; and Zou, J. 2023.
\newblock FrugalGPT.
\newblock \emph{arXiv:2305.05176}.

\bibitem[Dietterich(2000)]{dietterich2000ensemble}
Dietterich, T.~G. 2000.
\newblock Ensemble Methods in Machine Learning.
\newblock In \emph{Multiple Classifier Systems (MCS 2000)}, Lecture Notes in
Computer Science vol.\ 1857, 1--15. Springer.

\bibitem[Deng et~al.(2009)]{deng2009imagenet}
Deng, J.; Dong, W.; Socher, R.; Li, L.-J.; Li, K.; and Fei-Fei, L. 2009.
\newblock ImageNet: A Large-Scale Hierarchical Image Database.
\newblock In \emph{CVPR}, 248--255.

\bibitem[Gal and Ghahramani(2016)]{gal2016dropout}
Gal, Y.; and Ghahramani, Z. 2016.
\newblock Dropout as a Bayesian Approximation.
\newblock In \emph{ICML}, 1050--1059.

\bibitem[Geifman and El-Yaniv(2017)]{geifman2017selective}
Geifman, Y.; and El-Yaniv, R. 2017.
\newblock Selective Classification for Deep Neural Networks.
\newblock In \emph{NeurIPS}.

\bibitem[Groot and Valdenegro-Toro(2024)]{groot2024verbalized}
Groot, T.; and Valdenegro-Toro, M. 2024.
\newblock Overconfidence is Key: Verbalized Uncertainty Evaluation in Large
Language and Vision-Language Models.
\newblock In \emph{Proc.\ 4th Workshop on Trustworthy NLP (TrustNLP @ NAACL)}, 145--171.

\bibitem[Guo et~al.(2017)]{guo2017calibration}
Guo, C.; Pleiss, G.; Sun, Y.; and Weinberger, K.~Q. 2017.
\newblock On Calibration of Modern Neural Networks.
\newblock In \emph{ICML}, 1321--1330.

\bibitem[Hendrycks and Gimpel(2017)]{hendrycks2017baseline}
Hendrycks, D.; and Gimpel, K. 2017.
\newblock A Baseline for Detecting Misclassified and Out-of-Distribution Examples.
\newblock In \emph{ICLR}.

\bibitem[Kendall and Gal(2017)]{kendall2017uncertainties}
Kendall, A.; and Gal, Y. 2017.
\newblock What Uncertainties Do We Need in Bayesian Deep Learning for Computer Vision?
\newblock In \emph{NeurIPS}.

\bibitem[Krizhevsky(2009)]{krizhevsky2009learning}
Krizhevsky, A. 2009.
\newblock Learning Multiple Layers of Features from Tiny Images.
\newblock Tech.\ report, University of Toronto.

\bibitem[Lakshminarayanan, Pritzel, and Blundell(2017)]{lakshminarayanan2017simple}
Lakshminarayanan, B.; Pritzel, A.; and Blundell, C. 2017.
\newblock Simple and Scalable Predictive Uncertainty Estimation Using Deep Ensembles.
\newblock In \emph{NeurIPS}.

\bibitem[Lin(1991)]{lin1991divergence}
Lin, J. 1991.
\newblock Divergence Measures Based on the Shannon Entropy.
\newblock \emph{IEEE Trans.\ Information Theory}, 37(1): 145--151.

\bibitem[Naeini, Cooper, and Hauskrecht(2015)]{naeini2015obtaining}
Naeini, M.~P.; Cooper, G.~F.; and Hauskrecht, M. 2015.
\newblock Obtaining Well Calibrated Probabilities Using Bayesian Binning.
\newblock In \emph{AAAI}.

\bibitem[Nie, Zhou, and Bansal(2020)]{nie2020chaosnli}
Nie, Y.; Zhou, X.; and Bansal, M. 2020.
\newblock What Can We Learn from Collective Human Opinions on Natural Language
Inference Data?
\newblock In \emph{EMNLP}.

\bibitem[Northcutt, Athalye, and Mueller(2021)]{northcutt2021pervasive}
Northcutt, C.~G.; Athalye, A.; and Mueller, J. 2021.
\newblock Pervasive Label Errors in Test Sets Destabilize ML Benchmarks.
\newblock In \emph{NeurIPS Datasets and Benchmarks}.

\bibitem[Ong et~al.(2024)]{ong2024routellm}
Ong, I.; Almahairi, A.; Wu, V.; Chiang, W.-L.; Wu, T.; Gonzalez, J.~E.;
Kadous, M.~W.; and Stoica, I. 2024.
\newblock RouteLLM: Learning to Route LLMs with Preference Data.
\newblock \emph{arXiv:2406.18665}.

\bibitem[Peterson et~al.(2019)]{peterson2019human}
Peterson, J.~C.; Battleday, R.~M.; Griffiths, T.~L.; and Russakovsky, O. 2019.
\newblock Human Uncertainty Makes Classification More Robust.
\newblock In \emph{ICCV}, 9617--9626.

\bibitem[Platt(1999)]{platt1999probabilistic}
Platt, J.~C. 1999.
\newblock Probabilistic Outputs for Support Vector Machines.
\newblock In \emph{Advances in Large Margin Classifiers}, 61--74.

\bibitem[Ratner et~al.(2017)]{ratner2017snorkel}
Ratner, A.; Bach, S.~H.; Ehrenberg, H.; Fries, J.; Wu, S.; and R\'e, C. 2017.
\newblock Snorkel: Rapid Training Data Creation with Weak Supervision.
\newblock \emph{PVLDB}, 11(3): 269--282.

\bibitem[Reimers and Gurevych(2019)]{reimers2019sentencebert}
Reimers, N.; and Gurevych, I. 2019.
\newblock Sentence-BERT.
\newblock In \emph{EMNLP}, 3982--3992.

\bibitem[Settles(2009)]{settles2009active}
Settles, B. 2009.
\newblock Active Learning Literature Survey.
\newblock Tech.\ report 1648, Univ.\ of Wisconsin--Madison.

\bibitem[Viola and Jones(2001)]{viola2001rapid}
Viola, P.; and Jones, M. 2001.
\newblock Rapid Object Detection Using a Boosted Cascade of Simple Features.
\newblock In \emph{CVPR}.

\bibitem[Wang et~al.(2024)]{wang2024qwen2vl}
Wang, P.; et~al. 2024.
\newblock Qwen2-VL.
\newblock \emph{arXiv:2409.12191}.

\bibitem[Zadrozny and Elkan(2002)]{zadrozny2002transforming}
Zadrozny, B.; and Elkan, C. 2002.
\newblock Transforming Classifier Scores into Accurate Multiclass Probability Estimates.
\newblock In \emph{KDD}, 694--699.

\bibitem[Zhang, Zhang, and Zheng(2024)]{zhang2024vluncertainty}
Zhang, R.; Zhang, H.; and Zheng, Z. 2024.
\newblock VL-Uncertainty: Detecting Hallucination in Large Vision-Language
Model via Uncertainty Estimation.
\newblock \emph{arXiv:2411.11919}.

\bibitem[Zheng et~al.(2023)]{zheng2023judging}
Zheng, L.; et~al. 2023.
\newblock Judging LLM-as-a-Judge with MT-Bench and Chatbot Arena.
\newblock In \emph{NeurIPS Datasets and Benchmarks}.

\end{thebibliography}

\end{document}